\section{Related Work}

\paragraph{Karpathy's Recipe.} In 2019, Andrej Karpathy wrote a widely cited guide for the training of neural networks where he discusses how training fails silently \citep{karpathy2019}. He provides manual habits to catch these silent failures, and the Stanford CS231n course notes add checking that the update-to-weight ratio is near 0.1\% \citep{cs231n}. However, his guide simply provides a checklist a human needs to follow by hand with no mention of any automated solution.

\paragraph{wandb.ai.} Weights \& Biases (wandb) is the industry standard for experiment tracking through a dashboard where one line of code in a project can log live gradient and weight histograms during training \citep{biewald2020}. Still, it is limited by only displaying the numbers and nothing more, since a human still needs to watch, understand what healthy numbers look like, and catch the anomaly.

\paragraph{wandb-callbacks.} Wandb-callbacks is a small community add-on package for wandb that includes a dead-ReLU warning callback, which by default warns when at least 80\% of a layer's units never activate \citep{groeger_wandbcallbacks}. However, it covers only one failure mode with a single signal and provides no diagnosis or halting.

\paragraph{DeepLocalize.} DeepLocalize, published at the 2021 International Conference on Software Engineering, watches the values a network produces at each layer during training and, when it finds a numerical error, traces the error back to the layer that caused it \citep{wardat2021}. However, it only reports faults that lead to numerical errors and does not suggest how to fix them \citep{wardat2022}.

\paragraph{UMLAUT.} UMLAUT, presented at the 2021 ACM Conference on Human Factors in Computing Systems, checks a deep learning program's structure and its model's behavior against debugging heuristics drawn from experts, then shows readable error messages next to the code and the model's output in a single interface \citep{schoop2021}. In a study with 15 participants, people using it found and fixed significantly more bugs than without it. Still, its heuristics fire on models with no known bugs, and the authors of DeepDiagnosis observed it warning of possible overfitting on 68 of 78 such MNIST models \citep{wardat2022}.

\paragraph{SageMaker Debugger.} Amazon SageMaker Debugger records tensors during training jobs on Amazon's cloud and checks them with built-in rules such as DeadRelu, VanishingGradient, ExplodingTensor, UnchangedTensor, and WeightUpdateRatio \citep{aws_rules}. Any rule can be paired with an action that stops the training job when the rule fires \citep{aws_actions}, so detection and halting are already automated. However, its dead unit, gradient, and unchanged tensor rules compare a signal against a fixed default threshold. DeadRelu, for example, fires when more than 50\% of a layer's ReLUs are active less than 1\% of the time, and a user training a different network has to pick a new threshold by hand. Its SimilarAcrossRuns rule compares a job with one other completed job rather than with the range seen across many healthy runs. Amazon's documentation also states that Debugger is no longer open to new customers.

\paragraph{AutoTrainer.} AutoTrainer is a training monitor and automatic repair system published at the 2021 International Conference on Software Engineering \citep{zhang2021}. During training, it checks gradients, weights, and the loss and accuracy curves against preset thresholds to detect five problems: vanishing gradients, exploding gradients, dying ReLU, oscillating loss, and slow convergence. When it finds one, it retrains the model with built-in fixes such as replacing the activation function, adding batch normalization, or changing the weight initializer. However, it does not report which layer is at fault \citep{wardat2022}, and its default thresholds are the same for every network. A later study of model repair on 9 real and 21 artificial faulty models found that a random search did as well as or better than AutoTrainer and two hyperparameter tuning methods, and that all of them failed on larger models \citep{kim2023}.

\paragraph{DeepDiagnosis.} DeepDiagnosis, published at the 2022 International Conference on Software Engineering, checks a model at the end of every training epoch for eight error symptoms, including dead nodes, unchanged weights, saturated activations, and vanishing or exploding gradients \citep{wardat2022}. When it finds a symptom, it reports the layer, maps the symptom to a likely root cause, and suggests a fix, which is close to the layer, cause, and suggested check that LHD reports. However, its checks also rely on fixed thresholds, such as counting a layer as having dead nodes when more than 70\% of its ReLU units are inactive. On 188 models from the AutoTrainer benchmark with no known bugs, it still reported symptoms, including dead nodes in 21 of the 35 CIFAR-10 models \citep{wardat2022}.

\paragraph{ART.} Actually Robust Training is a Python library inspired by Karpathy's recipe that adds validation checks to training pipelines \citep{chwilczynski2024}. While this is a step up, it only imposes process steps at development time rather than live fault detection from internal signals during training.

\medskip

To conclude, tools that detect faults during training, and even stop or repair the run, already exist. However, their dead unit and gradient checks compare each signal against a threshold chosen in advance, such as 50\%, 70\%, or 80\% of a layer's units being dead, even though the share of units that die during healthy training differs between networks and between layers. A threshold that suits one network can therefore miss real faults in another or flag its healthy runs, as DeepDiagnosis's dead node check did on CIFAR-10 models with no known bugs. LHD instead measures a separate healthy range for every layer, signal, and epoch from healthy runs of the network being trained. Furthermore, DeepDiagnosis is evaluated by how many buggy models it flags and how quickly, not by how many epochs earlier its signals catch a fault than the accuracy curve or by how often it flags fresh healthy runs of the same network, and those are the two measurements this work reports.

